\chapter{Evaluation}
\label{chap:eval}

\scaffold{Chapter intro: four questions---fuzzing throughput, which violations are
found and how fast, which are Arm-specific, and whether they reproduce on commodity
hardware---plus end-to-end exploitation. Draw from uasc2027 ``Evaluation'' and the
arm-disclosure measurements.}

\section{Experimental setup}
\label{sec:eval-setup}
\scaffold{Neoverse~N3 (Armv9.2-A, QARMA3, 48-bit VAs, FEAT\_PAuth2/FPAC; GCP
c4a bare-metal; confirm TBI0/TBI1 and CONSTPACFIELD); Google Tensor~G3 (Pixel~8),
leak pinned to a Cortex-A715 core; Intel Alder Lake as a non-Arm baseline. Metric
definitions: throughput and time-to-detection (TTD, on the N3).}

\section{Fuzzing throughput}
\label{sec:eval-throughput}
\scaffold{Local versus remote orchestration throughput table (test cases/hour, inputs
/test case, executions/second). Numbers to be filled.}

\section{Contract violations}
\label{sec:eval-violations}
\scaffold{The violations table: Spectre-v1, Spectre-v2 (L1D, standard; and BTB, the
case current priming misses), Spectre-v4, the canonicality leak (cache/TLB and BTB),
the AUT* sign-bit leak, and the unexplained L1D IP-prefetcher effect (Alder Lake
only). TTD on the N3; reproduction on N3 / Pixel~8 (A715) / Alder Lake.}

\section{Localization cost}
\label{sec:eval-localize}
\scaffold{Hardware executions to localize a boundary: linear scan ($k$) versus binary
search ($\lceil\log_2(k+1)\rceil$), tying back to \autoref{chap:priming} and the
proofs in \autoref{appendix:localization}.}

\section{End-to-end exploitation}
\label{sec:eval-exploit}
\scaffold{Cross-boundary recovery (cross-process BTB: 100\% over the PMU, $\approx$65
bit/s; guest-to-host BTB: 100\%, $\approx$130 bit/s; EL2$\to$EL0 dTLB: 100\%,
$\approx$1230 bit/s). A ROP chain built entirely from OpenSSL's own gadgets recovers
2044/2048 bits ($99.8\%$) of an RSA-2048 private key. Blind stack-canary recovery via
three windows (v1 20/20, v4 18/20, rsb 19/20). Reconcile the ``musl'' ROP chain from
the abstract against the current artifacts (open item).}

\section{Covert-channel analysis}
\label{sec:eval-covert}
\scaffold{Capacity and error of the channel; the context-switch wall that forces the
cross-domain variant onto the BTB, trading bandwidth for persistence.}
